\documentclass[UTF8,a4paper,11pt]{ctexart}

\usepackage[margin=2.2cm]{geometry}
\usepackage{hyperref}
\usepackage{booktabs}
\usepackage{tabularx}
\usepackage{array}
\usepackage{longtable}
\usepackage{enumitem}
\usepackage{xcolor}
\usepackage{colortbl}
\usepackage{ragged2e}
\usepackage{amsmath}

\hypersetup{
  colorlinks=true,
  linkcolor=blue!50!black,
  urlcolor=teal!60!black,
  citecolor=blue!50!black
}
\urlstyle{tt}

\definecolor{mmBlue}{HTML}{123B5D}
\definecolor{mmTeal}{HTML}{1E6F78}
\definecolor{mmGold}{HTML}{C98B2E}
\definecolor{mmGray}{HTML}{5B6573}
\definecolor{mmLite}{HTML}{EEF3F7}
\definecolor{mmMustBg}{HTML}{FDECEC}
\definecolor{mmLaterBg}{HTML}{FFF6E5}
\definecolor{mmSkipBg}{HTML}{EDF2F7}

\newcolumntype{Y}{>{\RaggedRight\arraybackslash}X}
\newcolumntype{P}[1]{>{\RaggedRight\arraybackslash}p{#1}}
\newcommand{\kpath}[1]{\path{#1}}
\newcommand{\musttag}{\colorbox{mmMustBg}{\textbf{必须掌握}}}
\newcommand{\latertag}{\colorbox{mmLaterBg}{\textbf{后续再看}}}
\newcommand{\skiptag}{\colorbox{mmSkipBg}{\textbf{暂时跳过}}}

\setlist[itemize]{leftmargin=2em,itemsep=0.25em,topsep=0.35em}
\setlist[enumerate]{leftmargin=2em,itemsep=0.25em,topsep=0.35em}

\title{\textbf{Linux \texttt{mm} 初学者学习文档}}
\author{基于仓库内源码树 \texttt{Linux 7.0.0-rc7} 整理}
\date{生成日期：2026-04-12}

\begin{document}

\maketitle
\tableofcontents
\newpage

\section{文档定位}

这份文档面向“刚开始系统学习 Linux 内存管理（\texttt{mm}）”的读者，目标不是覆盖
\texttt{mm/} 目录下所有子系统，而是提供一条对初学者成本最低、收益最高的阅读路径。

\paragraph{本文的使用方式}
\begin{itemize}
  \item 先按本文给出的顺序建立整体地图，再去读代码。
  \item 每一块都用三种优先级标记：
  \begin{itemize}
    \item \musttag：第一轮必须看懂，否则后面的代码会断层。
    \item \latertag：建议在主路径打通后再补。
    \item \skiptag：先不要投入时间，避免刚入门就被支线拖住。
  \end{itemize}
  \item 这份文档以仓库内源码树为准。外部链接在 2026-04-12 检查过，但在线文档可能比本地
  \texttt{7.0.0-rc7} 略有漂移；读实现细节时优先看本地源码。
\end{itemize}

\paragraph{阅读总原则}
\begin{enumerate}
  \item 先弄清“对象和层次”，再看算法细节。
  \item 先打通主路径，再进入优化路径和异常路径。
  \item 先选一套架构 fault 路径，推荐只看 \texttt{x86}。
  \item 不要一上来平铺阅读 \texttt{mm/*.c}；这样最容易迷路。
\end{enumerate}

\section{整体路线总表}

\begin{longtable}{P{0.16\textwidth}P{0.22\textwidth}P{0.23\textwidth}P{0.31\textwidth}}
\toprule
主题块 & 先看文档 & 先看代码 & 第一轮目标 \\
\midrule
\endhead
总览与术语 &
\kpath{Documentation/mm/index.rst} \newline
\kpath{Documentation/admin-guide/mm/concepts.rst} &
\kpath{include/linux/mm_types.h} \newline
\kpath{include/linux/mmzone.h} \newline
\kpath{include/linux/gfp_types.h} &
建立 node/zone/folio/VMA/page table/reclaim 的全局地图。 \\
\addlinespace
物理内存与页分配 &
\kpath{Documentation/mm/physical_memory.rst} &
\kpath{mm/page_alloc.c} \newline
\kpath{mm/compaction.c} &
理解 page allocator 在分配、回收、碎片化之间如何联动。 \\
\addlinespace
地址空间、页表与缺页 &
\kpath{Documentation/mm/process_addrs.rst} \newline
\kpath{Documentation/mm/page_tables.rst} &
\kpath{mm/mmap.c} \newline
\kpath{mm/vma.c} \newline
\kpath{mm/memory.c} \newline
\kpath{arch/x86/mm/fault.c} &
打通 \texttt{mmap -> page fault -> anon/file fault -> COW} 主路径。 \\
\addlinespace
页缓存与文件映射 &
\kpath{Documentation/mm/page_cache.rst} &
\kpath{mm/filemap.c} \newline
\kpath{mm/readahead.c} \newline
\kpath{mm/workingset.c} &
理解文件读写、\texttt{mmap} 文件页 fault、readahead 和 refault。 \\
\addlinespace
回收、交换与 OOM &
\kpath{Documentation/admin-guide/mm/concepts.rst} &
\kpath{mm/vmscan.c} \newline
\kpath{mm/rmap.c} \newline
\kpath{mm/swap\_state.c} \newline
\kpath{mm/swapfile.c} \newline
\kpath{mm/oom\_kill.c} &
理解为什么能回收、如何回收、何时会 swap/OOM。 \\
\addlinespace
进阶主题 &
\kpath{Documentation/mm/transhuge.rst} \newline
\kpath{Documentation/mm/multigen\_lru.rst} \newline
\kpath{Documentation/mm/numa.rst} &
\kpath{mm/huge\_memory.c} \newline
\kpath{mm/khugepaged.c} \newline
\kpath{mm/memcontrol.c} &
在主线走通后再进入 THP、memcg、NUMA、MGLRU。 \\
\bottomrule
\end{longtable}

\section{第 0 块：先建立脑图，不要急着看所有子系统}

\subsection*{这一步必须掌握的对象}
\begin{itemize}
  \item \musttag \textbf{物理内存层：} node、zone、page frame、folio。
  \item \musttag \textbf{虚拟地址空间层：} \texttt{mm\_struct}、\texttt{vm\_area\_struct}、page table。
  \item \musttag \textbf{文件与缓存层：} \texttt{address\_space}、page cache、readahead、workingset。
  \item \musttag \textbf{压力处理层：} reclaim、swap、OOM、compaction。
  \item \latertag \textbf{控制与策略层：} THP、memcg、NUMA policy、MGLRU、DAMON、HMM。
\end{itemize}

\subsection*{最先应该记住的结构体}
\begin{itemize}
  \item \musttag \kpath{include/linux/mm_types.h} 里的 \texttt{struct page}、\texttt{struct folio}、
  \texttt{struct vm\_area\_struct}、\texttt{struct mm\_struct}。
  \item \musttag \kpath{include/linux/mmzone.h} 里的 \texttt{struct zone}、
  \texttt{struct per\_cpu\_pages}、\texttt{struct pglist\_data}。
  \item \musttag \kpath{include/linux/gfp_types.h} 里的 \texttt{gfp\_t} 和各种 \texttt{\_\_GFP\_*} 语义。
\end{itemize}

\subsection*{现在先不要深挖的内容}
\begin{itemize}
  \item \skiptag \texttt{DAMON}、\texttt{HMM}、\texttt{secretmem}、\texttt{hwpoison}、\texttt{userfaultfd}。
  \item \skiptag \texttt{CMA}、\texttt{memory\_hotplug}、\texttt{balloon}、\texttt{memory-tiers}。
  \item \skiptag \texttt{nommu} 和大部分 \texttt{highmem} 细节，除非你明确要做 32 位内核学习。
\end{itemize}

\subsection*{外部参考}
\begin{itemize}
  \item 官方总览：\url{https://docs.kernel.org/mm/index.html}
  \item 概念总览：\url{https://docs.kernel.org/admin-guide/mm/concepts.html}
  \item 历史技术书目录（适合查章节，不适合直接对号入座当前实现）：
  \url{https://www.kernel.org/doc/gorman/html/understand/understand001.html}
\end{itemize}

\section{第 1 块：物理内存、页分配与分配失败后的联动}

\subsection*{为什么从这里开始}
如果你不知道 page allocator 在做什么，那么后面看到 page fault、page cache、slab、swap
时都会停留在“它又去拿了一页”的表面层。Linux \texttt{mm} 的很多路径最终都会落到
\texttt{alloc\_pages()} 或其慢路径上。

\subsection*{本块重点知识点}
\begin{itemize}
  \item \musttag \textbf{node / zone / zonelist。} 分配不是在“全局内存池”里随便拿，而是按 node、
  zone 和 fallback 顺序选择。
  \item \musttag \textbf{order 与 buddy allocator。} 理解 order-0、order-\(n\)、split、merge
  和外部碎片化。
  \item \musttag \textbf{per-cpu pagesets。} 绝大多数 order-0 分配并不直接碰全局 buddy。
  \item \musttag \textbf{GFP 标志。} \texttt{GFP\_KERNEL}、\texttt{\_\_GFP\_DIRECT\_RECLAIM}、
  \texttt{\_\_GFP\_KSWAPD\_RECLAIM}、\texttt{\_\_GFP\_HIGH}、\texttt{\_\_GFP\_THISNODE}
  到底影响什么。
  \item \musttag \textbf{快路径和慢路径。} 先成功从 freelist/pcp 拿页；失败后才进入 direct reclaim、
  compaction、OOM 相关慢路径。
  \item \latertag \textbf{watermark、boost、reserve。} 第一轮要知道它们存在并影响回收时机，
  不必一开始背所有细节。
  \item \latertag \textbf{migratetype 与 pageblock。} 它们和 compaction、CMA、碎片化治理强相关。
  \item \skiptag \textbf{memblock、sparsemem、vmemmap dedup。} 这些更适合作为第二轮专题。
\end{itemize}

\subsection*{第一轮源码入口}
\begin{itemize}
  \item \musttag \kpath{mm/page_alloc.c}：
  \begin{itemize}
    \item \texttt{\_\_alloc\_pages\_noprof()}
    \item \texttt{get\_page\_from\_freelist()}
    \item \texttt{\_\_alloc\_pages\_slowpath()}
    \item \texttt{rmqueue()}
    \item \texttt{\_\_free\_one\_page()}
  \end{itemize}
  \item \musttag \kpath{include/linux/mmzone.h}：先看 \texttt{enum zone\_type}、
  \texttt{struct zone}、\texttt{struct per\_cpu\_pages}、\texttt{struct pglist\_data}。
  \item \musttag \kpath{include/linux/gfp_types.h}：先建立 GFP 语义，不要死记组合宏。
  \item \latertag \kpath{mm/compaction.c}：先知道高阶分配失败后为什么会牵出 compaction。
\end{itemize}

\subsection*{建议阅读顺序}
\begin{enumerate}
  \item \kpath{Documentation/mm/physical_memory.rst}
  \item \kpath{include/linux/mmzone.h}
  \item \kpath{include/linux/gfp_types.h}
  \item \kpath{mm/page_alloc.c}
  \item 再回来看一次 \texttt{page\_alloc.c} 里的 slowpath
\end{enumerate}

\subsection*{对应外部资料}
\begin{itemize}
  \item 官方文档：\url{https://docs.kernel.org/mm/physical_memory.html}
  \item Mel Gorman，\textit{Describing Physical Memory}：
  \url{https://www.kernel.org/doc/gorman/html/understand/understand005.html}
  \item Mel Gorman，\textit{Physical Page Allocation}：
  \url{https://www.kernel.org/doc/gorman/html/understand/understand009.html}
  \item Linux-MM wiki，\textit{PageAllocation}：
  \url{https://linux-mm.org/PageAllocation}
\end{itemize}

\paragraph{如何使用这些外部资料}
\begin{itemize}
  \item 官方文档用于对照当前对象和术语。
  \item Gorman 的书适合建立 buddy、zone、per-cpu pages 的基本直觉，但字段和函数名有时代差。
  \item Linux-MM wiki 适合把页分配器放回设计背景里看，不要直接拿它解释当前代码分支。
\end{itemize}

\section{第 2 块：进程地址空间、VMA、页表与缺页异常}

\subsection*{为什么这是整条主线的核心}
对新手来说，Linux \texttt{mm} 最重要的闭环不是“某个算法”，而是：

\begin{center}
\texttt{mmap/brk} \(\rightarrow\) 建立 VMA \(\rightarrow\) CPU 访问地址 \(\rightarrow\) page fault
\(\rightarrow\) \texttt{handle\_mm\_fault()} \(\rightarrow\) 匿名页 / 文件页 / COW / swap fault
\end{center}

只要这条路径走通，后面看 \texttt{rmap}、reclaim、swap、THP 都会容易很多。

\subsection*{本块重点知识点}
\begin{itemize}
  \item \musttag \textbf{\texttt{mm\_struct} 和 \texttt{vm\_area\_struct}。}
  一个进程地址空间有哪些 VMA，它们如何被维护。
  \item \musttag \textbf{maple tree。} 当前主线不再是老资料里常说的 VMA 红黑树主视角。
  \item \musttag \textbf{VMA 锁语义。} \texttt{mmap\_lock}、VMA lock、rmap lock 是不同层次的问题。
  \item \musttag \textbf{五级页表抽象。} PGD/P4D/PUD/PMD/PTE 的层次、folding 和 huge mapping。
  \item \musttag \textbf{匿名缺页、文件缺页、写时复制（COW）。}
  \item \musttag \textbf{fault 路径的入口。} 对 x86 来说，从
  \kpath{arch/x86/mm/fault.c} 进 \texttt{handle\_mm\_fault()}。
  \item \latertag \textbf{NUMA fault、per-VMA lock 优化、huge fault。} 第一轮知道它们在主路径里出现即可。
  \item \latertag \textbf{mprotect/mremap/userfaultfd/GUP/mmu notifier。}
  它们很重要，但不应先于基础缺页路径。
\end{itemize}

\subsection*{第一轮必看源码}
\begin{itemize}
  \item \musttag \kpath{include/linux/mm_types.h}：\texttt{struct vm\_area\_struct}、
  \texttt{struct mm\_struct}。
  \item \musttag \kpath{mm/mmap.c}：\texttt{do\_mmap()}。
  \item \musttag \kpath{mm/vma.c}：\texttt{mmap\_region()}、\texttt{vma\_merge\_new\_range()}、
  \texttt{do\_vmi\_munmap()}。
  \item \musttag \kpath{mm/memory.c}：\texttt{handle\_mm\_fault()}、
  \texttt{\_\_handle\_mm\_fault()}、\texttt{do\_anonymous\_page()}、
  \texttt{do\_wp\_page()}、\texttt{do\_swap\_page()}、\texttt{finish\_fault()}。
  \item \musttag \kpath{arch/x86/mm/fault.c}：\texttt{do\_user\_addr\_fault()}，
  只选一个架构先看。
\end{itemize}

\subsection*{建议阅读顺序}
\begin{enumerate}
  \item \kpath{Documentation/mm/process_addrs.rst}
  \item \kpath{Documentation/mm/page_tables.rst}
  \item \kpath{include/linux/mm_types.h}
  \item \kpath{mm/mmap.c} 与 \kpath{mm/vma.c}
  \item \kpath{arch/x86/mm/fault.c}
  \item \kpath{mm/memory.c}
\end{enumerate}

\subsection*{第一轮可以先不看的内容}
\begin{itemize}
  \item \skiptag \texttt{gup.c} 和 \texttt{process\_vm\_access.c}
  \item \skiptag \texttt{userfaultfd.c}
  \item \skiptag \texttt{mmu\_notifier.c}
  \item \skiptag \texttt{mseal.c}、\texttt{secretmem.c} 等专项功能
\end{itemize}

\subsection*{对应外部资料}
\begin{itemize}
  \item 官方文档，\textit{Process Addresses}：
  \url{https://docs.kernel.org/mm/process_addrs.html}
  \item 官方文档，\textit{Page Tables}：
  \url{https://docs.kernel.org/mm/page_tables.html}
  \item Mel Gorman，\textit{Page Table Management}：
  \url{https://www.kernel.org/doc/gorman/html/understand/understand006.html}
  \item Mel Gorman，\textit{Process Address Space}：
  \url{https://www.kernel.org/doc/gorman/html/understand/understand007.html}
  \item \textit{active\_mm} 背景说明（Linus 在 lore 上的原始解释）：
  \url{https://lore.kernel.org/lkml/Pine.LNX.4.10.9907301410280.752-100000@penguin.transmeta.com/}
\end{itemize}

\paragraph{特别提醒}
很多旧资料会说“VMA 放在红黑树里”，这对理解历史没问题，但对当前主线不够准确。
现在一定要把“\texttt{mm} 内部使用 maple tree 管理 VMA”记牢。

\section{第 3 块：页缓存、文件映射与文件页 fault}

\subsection*{为什么这一块很早就要看}
初学者往往把“缺页”误认为只处理匿名页。实际上文件映射、普通 \texttt{read()}、readahead、
cache refault、workingset 这些路径，决定了你能不能真正理解 Linux 如何把文件系统和内存管理
接起来。

\subsection*{本块重点知识点}
\begin{itemize}
  \item \musttag \textbf{page cache 的对象模型。} 文件页现在主要以 \texttt{folio} 为单位管理。
  \item \musttag \textbf{\texttt{address\_space}。} 文件页不是凭空存在，它们挂在 mapping 上。
  \item \musttag \textbf{文件页 fault。} \texttt{filemap\_fault()} 与匿名缺页是两条不同但会汇合的路径。
  \item \musttag \textbf{readahead。} 同步/异步 readahead 与实际 cache miss 路径的关系。
  \item \musttag \textbf{workingset / refault。} 这是当前 reclaim 与 page cache 反馈的关键组成部分。
  \item \latertag \textbf{truncate、writeback、dirty throttling。} 很重要，但第一轮只需建立接口感。
  \item \latertag \textbf{tmpfs/shmem。} 它把“匿名”和“文件页”的边界进一步揉在一起，第二轮再细看。
\end{itemize}

\subsection*{第一轮必看源码}
\begin{itemize}
  \item \musttag \kpath{mm/filemap.c}：\texttt{filemap\_fault()}、\texttt{filemap\_read()}。
  \item \musttag \kpath{mm/readahead.c}：\texttt{page\_cache\_ra\_unbounded()}。
  \item \musttag \kpath{mm/workingset.c}：
  \texttt{workingset\_eviction()}、\texttt{workingset\_refault()}、
  \texttt{workingset\_activation()}。
  \item \latertag \kpath{mm/shmem.c}：理解 tmpfs/shmem 时再细看。
\end{itemize}

\subsection*{建议阅读顺序}
\begin{enumerate}
  \item \kpath{Documentation/mm/page_cache.rst}
  \item \kpath{include/linux/pagemap.h}
  \item \kpath{mm/filemap.c}
  \item \kpath{mm/readahead.c}
  \item \kpath{mm/workingset.c}
\end{enumerate}

\subsection*{这块先不用深挖的内容}
\begin{itemize}
  \item \skiptag \kpath{mm/page-writeback.c}
  \item \skiptag \kpath{mm/mapping_dirty_helpers.c}
  \item \skiptag \kpath{mm/fadvise.c}、\kpath{mm/msync.c}
\end{itemize}

\subsection*{对应外部资料}
\begin{itemize}
  \item 官方文档，\textit{Page Cache}：
  \url{https://docs.kernel.org/next/mm/page_cache.html}
  \item 官方文档，\textit{Memory Management APIs / Filemap}：
  \url{https://docs.kernel.org/core-api/mm-api.html}
  \item Mel Gorman，\textit{Page Frame Reclamation} 中的 page cache 部分：
  \url{https://www.kernel.org/doc/gorman/html/understand/understand013.html}
\end{itemize}

\paragraph{特别提醒}
如果你看的资料大量使用 “page” 作为页缓存核心单位，不必惊慌。当前主线里很多接口已经 folio 化，
但核心问题还是那几个：页在不在缓存里、是不是 uptodate、要不要触发 readahead、
refault 能不能说明当前 working set 被挤掉了。

\section{第 4 块：回收、反向映射、交换与 OOM}

\subsection*{为什么这一块不能拖到最后}
一旦分配失败、内存变紧，Linux \texttt{mm} 的主线就会迅速转向 reclaim、rmap、swap、OOM。
这部分如果不懂，你就解释不了：

\begin{itemize}
  \item 为什么某些页能被回收，某些不能；
  \item 为什么 reclaim 需要知道一个 folio 被哪些页表映射；
  \item 为什么有时先 direct reclaim，有时唤醒 \texttt{kswapd}；
  \item 为什么最后会走到 OOM killer。
\end{itemize}

\subsection*{本块重点知识点}
\begin{itemize}
  \item \musttag \textbf{reclaimable 与 unreclaimable 的区别。}
  \item \musttag \textbf{file/anon 两类页在 reclaim 中的不同命运。}
  \item \musttag \textbf{rmap 的目的。} 不是“多维护一个映射”而已，而是 reclaim/COW/migration 都要靠它反查。
  \item \musttag \textbf{try\_to\_unmap。} 页要被回收前，往往先得把页表映射拆掉。
  \item \musttag \textbf{kswapd 与 direct reclaim。}
  \item \musttag \textbf{swap cache 与 swapin readahead。}
  \item \musttag \textbf{OOM 的触发位置与分工。}
  \item \latertag \textbf{workingset feedback、MGLRU、refault distance。}
  \item \latertag \textbf{zswap/zsmalloc。} 它们是 swap 的优化层，不是第一层机制。
\end{itemize}

\subsection*{第一轮必看源码}
\begin{itemize}
  \item \musttag \kpath{mm/rmap.c}：
  \texttt{folio\_add\_new\_anon\_rmap()}、\texttt{rmap\_walk()}、
  \texttt{try\_to\_unmap()}、\texttt{folio\_referenced()}。
  \item \musttag \kpath{mm/vmscan.c}：
  \texttt{try\_to\_free\_pages()}、\texttt{shrink\_node()}、
  \texttt{balance\_pgdat()}、\texttt{wakeup\_kswapd()}。
  \item \musttag \kpath{mm/swap\_state.c}：
  \texttt{read\_swap\_cache\_async()}、\texttt{swapin\_readahead()}。
  \item \musttag \kpath{mm/swapfile.c}：先建立 \texttt{swap\_info\_struct} 和 swap area 的概念。
  \item \musttag \kpath{mm/oom\_kill.c}：
  \texttt{out\_of\_memory()}、\texttt{select\_bad\_process()}、
  \texttt{oom\_kill\_process()}。
\end{itemize}

\subsection*{建议阅读顺序}
\begin{enumerate}
  \item 先重读 \kpath{Documentation/admin-guide/mm/concepts.rst} 里的 reclaim/compaction/OOM 部分
  \item \kpath{mm/rmap.c}
  \item \kpath{mm/vmscan.c}
  \item \kpath{mm/swap\_state.c} 与 \kpath{mm/swapfile.c}
  \item \kpath{mm/oom\_kill.c}
\end{enumerate}

\subsection*{第一轮可以暂缓}
\begin{itemize}
  \item \skiptag \kpath{mm/zswap.c}
  \item \skiptag \kpath{mm/zsmalloc.c}
  \item \skiptag \kpath{mm/page_owner.c}
  \item \skiptag \kpath{mm/page_table_check.c}
\end{itemize}

\subsection*{对应外部资料}
\begin{itemize}
  \item 官方概念文档（reclaim/OOM 总览）：
  \url{https://docs.kernel.org/admin-guide/mm/concepts.html}
  \item 官方文档，\textit{Multi-Gen LRU}：
  \url{https://docs.kernel.org/admin-guide/mm/multigen_lru.html}
  \item Mel Gorman，\textit{Page Frame Reclamation}：
  \url{https://www.kernel.org/doc/gorman/html/understand/understand013.html}
  \item Mel Gorman，\textit{Swap Management}：
  \url{https://www.kernel.org/doc/gorman/html/understand/understand014.html}
  \item Linux-MM wiki，\textit{PageReplacementDesign}：
  \url{https://linux-mm.org/PageReplacementDesign}
\end{itemize}

\paragraph{如何读这块}
第一轮不要试图把 \kpath{mm/vmscan.c} 每个分支一次性吃掉。先抓四个问题：

\begin{enumerate}
  \item 谁在发起 reclaim？
  \item 这次优先回收 file 还是 anon？
  \item 页在回收前为什么需要 rmap/try-to-unmap？
  \item reclaim 失败后什么时候会走到 OOM？
\end{enumerate}

\section{第 5 块：进阶主题，应该什么时候补}

\subsection*{这一块不是“不重要”，而是“不应抢跑”}
THP、memcg、NUMA、MGLRU、DAMON、HMM 都是现代 Linux \texttt{mm} 的关键主题，但它们都建立在
前面几块已经打通的前提之上。对初学者来说，最常见的问题不是“这些特性太难”，而是
“在主线还没通的时候太早进去”。

\subsection*{推荐顺序}
\begin{enumerate}
  \item \musttag THP：先看 \texttt{PMD}-sized huge page fault、collapse、khugepaged。
  \item \latertag MGLRU：在你已经理解传统 reclaim、refault、workingset 之后再进入。
  \item \latertag memcg：在你已经知道全局回收是怎么回事后，再看 cgroup 维度的回收。
  \item \latertag NUMA：先掌握 node/zone，再进 mempolicy 和自动 NUMA balancing。
  \item \skiptag DAMON/HMM/memory tiers：这些更像专题，不适合第一轮。
\end{enumerate}

\subsection*{推荐入口}
\begin{itemize}
  \item THP：
  \kpath{Documentation/mm/transhuge.rst}，
  \kpath{Documentation/admin-guide/mm/transhuge.rst}，
  \kpath{mm/huge\_memory.c}，
  \kpath{mm/khugepaged.c}
  \item MGLRU：
  \kpath{Documentation/mm/multigen\_lru.rst}，
  \kpath{Documentation/admin-guide/mm/multigen\_lru.rst}
  \item memcg：
  \kpath{mm/memcontrol.c}
  \item NUMA：
  \kpath{Documentation/mm/numa.rst}，
  \kpath{Documentation/admin-guide/mm/numa\_memory\_policy.rst}，
  \kpath{mm/mempolicy.c}
\end{itemize}

\subsection*{对应外部资料}
\begin{itemize}
  \item THP 官方文档：
  \url{https://docs.kernel.org/admin-guide/mm/transhuge.html}
  \item MGLRU 官方文档：
  \url{https://docs.kernel.org/admin-guide/mm/multigen_lru.html}
  \item 进程页表检查工具：
  \url{https://docs.kernel.org/admin-guide/mm/pagemap.html}
\end{itemize}

\section{初学者第一轮明确可以不看的文件}

下面这些不是“不重要”，而是“对第一轮收益不高”：

\begin{itemize}
  \item \skiptag \kpath{mm/damon/}
  \item \skiptag \kpath{mm/hmm.c}
  \item \skiptag \kpath{mm/userfaultfd.c}
  \item \skiptag \kpath{mm/secretmem.c}
  \item \skiptag \kpath{mm/hwpoison*}
  \item \skiptag \kpath{mm/cma*}
  \item \skiptag \kpath{mm/memory_hotplug.c}
  \item \skiptag \kpath{mm/page_owner.c}、\kpath{mm/debug*.c}、\kpath{mm/*test*.c}
  \item \skiptag 同时比较 \texttt{x86/arm64/riscv} 三套 fault 路径
\end{itemize}

\section{推荐的四周学习安排}

\subsection*{第 1 周：地图与基础对象}
\begin{itemize}
  \item 目标：建立全局术语地图。
  \item 阅读：\kpath{Documentation/admin-guide/mm/concepts.rst}、
  \kpath{Documentation/mm/physical_memory.rst}、
  \kpath{Documentation/mm/page_tables.rst}、
  \kpath{Documentation/mm/process_addrs.rst}。
  \item 代码：\kpath{include/linux/mm_types.h}、
  \kpath{include/linux/mmzone.h}、
  \kpath{include/linux/gfp_types.h}。
\end{itemize}

\subsection*{第 2 周：分配与地址空间}
\begin{itemize}
  \item 目标：理解 \texttt{alloc\_pages()} 和 \texttt{mmap()} 是什么。
  \item 代码：\kpath{mm/page_alloc.c}、\kpath{mm/mmap.c}、\kpath{mm/vma.c}。
  \item 产出：能画出一个 VMA 被建立出来的大致流程。
\end{itemize}

\subsection*{第 3 周：缺页与文件映射}
\begin{itemize}
  \item 目标：打通 \texttt{page fault} 主路径。
  \item 代码：\kpath{arch/x86/mm/fault.c}、\kpath{mm/memory.c}、
  \kpath{mm/filemap.c}、\kpath{mm/readahead.c}。
  \item 产出：能解释 anon fault、file fault、COW 的区别。
\end{itemize}

\subsection*{第 4 周：回收与交换}
\begin{itemize}
  \item 目标：理解内存压力下系统为什么这样行动。
  \item 代码：\kpath{mm/rmap.c}、\kpath{mm/vmscan.c}、
  \kpath{mm/swap\_state.c}、\kpath{mm/swapfile.c}、
  \kpath{mm/oom\_kill.c}。
  \item 产出：能解释 direct reclaim、\texttt{kswapd}、swap、OOM 之间的关系。
\end{itemize}

\section{外部参考汇总}

\subsection*{官方与当前文档}
\begin{itemize}
  \item \url{https://docs.kernel.org/mm/index.html}
  \item \url{https://docs.kernel.org/admin-guide/mm/concepts.html}
  \item \url{https://docs.kernel.org/mm/physical_memory.html}
  \item \url{https://docs.kernel.org/mm/page_tables.html}
  \item \url{https://docs.kernel.org/mm/process_addrs.html}
  \item \url{https://docs.kernel.org/next/mm/page_cache.html}
  \item \url{https://docs.kernel.org/admin-guide/mm/transhuge.html}
  \item \url{https://docs.kernel.org/admin-guide/mm/multigen_lru.html}
  \item \url{https://docs.kernel.org/admin-guide/mm/pagemap.html}
\end{itemize}

\subsection*{历史技术文档与长期参考}
\begin{itemize}
  \item Mel Gorman, \textit{Understanding the Linux Virtual Memory Manager}：
  \url{https://www.kernel.org/doc/gorman/html/understand/understand001.html}
  \item \textit{Describing Physical Memory}：
  \url{https://www.kernel.org/doc/gorman/html/understand/understand005.html}
  \item \textit{Page Table Management}：
  \url{https://www.kernel.org/doc/gorman/html/understand/understand006.html}
  \item \textit{Process Address Space}：
  \url{https://www.kernel.org/doc/gorman/html/understand/understand007.html}
  \item \textit{Physical Page Allocation}：
  \url{https://www.kernel.org/doc/gorman/html/understand/understand009.html}
  \item \textit{Page Frame Reclamation}：
  \url{https://www.kernel.org/doc/gorman/html/understand/understand013.html}
  \item \textit{Swap Management}：
  \url{https://www.kernel.org/doc/gorman/html/understand/understand014.html}
  \item Linux-MM wiki, \textit{PageAllocation}：
  \url{https://linux-mm.org/PageAllocation}
  \item Linux-MM wiki, \textit{PageReplacementDesign}：
  \url{https://linux-mm.org/PageReplacementDesign}
\end{itemize}

\paragraph{如何区分这些资料的用途}
\begin{itemize}
  \item 官方文档：适合对照当前主线术语和对象。
  \item Gorman：适合理解原理和主线结构，不适合直接对照今天的字段名、锁和具体函数。
  \item Linux-MM wiki：适合理解设计动机和历史背景，不适合当作当前实现说明书。
\end{itemize}

\section{结论}

对初学者来说，Linux \texttt{mm} 最重要的不是一上来知道所有子系统，而是先打通这五件事：

\begin{enumerate}
  \item 物理内存如何被组织；
  \item 地址空间如何被描述；
  \item 缺页如何被处理；
  \item 文件页如何进入和离开 page cache；
  \item 内存不够时 reclaim/swap/OOM 如何联动。
\end{enumerate}

只要这五块走通，再进入 THP、MGLRU、memcg、NUMA，就会是“在现有地图上扩展”，而不是
“重新学一套新系统”。

\end{document}
